\section{The crisis adoption wave}
\label{sec:crisis}

\subsection{Who adopts at the steady state}
\label{sec:ss_adoption}

At the calibrated steady state green durables are $5\%$ of the stock and the
quarterly switching flow is $0.26\%$, the flow that replaces green depreciation
(Table~\ref{tab:switch_prob_dist}, Appendix~\ref{app:taste}). Adoption is far from
uniform across households. Figure~\ref{fig:switch_heatmap} plots the switching
probability of the median discount-factor type over the wealth--productivity
grid: the choice \emph{policy function} rather than the mass-weighted pool, which mixes
three types whose switching rates differ by orders of magnitude and is best read as
the distribution in Table~\ref{tab:switch_prob_dist}. The probability that a brown incumbent
switches to green rises monotonically in both wealth and idiosyncratic
productivity. Wealthier households can finance the switching cost $\psi_g$ and more
productive households expect the higher future income that makes the durable worth
replacing, so adoption concentrates among wealthy, high-productivity households. The
surface also makes the rare-margin regime visible: away from that group the
probability sits far below one. The brown mass occupies the exponential tail of the
logit, where the semi-elasticity of switching is $1/\sigma_\varepsilon$ household by
household (Appendix~\ref{app:taste}, Table~\ref{tab:switch_prob_dist}).

\begin{figure}[H]
\centering
\figtitle{Switching probability over the wealth--productivity grid}
\includegraphics[width=0.72\textwidth]{output/fig_switch_heatmap.png}
\label{fig:switch_heatmap}
\fignotes{Probability of switching from a brown to a green durable, per quarter, for
the median discount-factor type, as a function of household wealth relative to the
mean $a/\bar a$ (horizontal, symmetric-log scale) and idiosyncratic productivity $e$
(vertical), at the baseline steady state. Colour on a log scale. A single type's
policy function is plotted rather than the mass-weighted pool: pooling the three
discount-factor types, whose switching rates differ by orders of magnitude, makes
the wealth gradient non-monotone even though each type's rule is monotone.}
\end{figure}

The wealth gradient is the marginal-switcher condition \eqref{eq:margin} seen
household by household. Figure~\ref{fig:payoff} evaluates one switch the way the
condition does: the upfront cost $\psi_g$ against the energy saving
$(\bar P^E_B-\bar P^E_G)\,\bar c^{\,G}_E$ accumulated over the durable's life,
discounted by interest and survival; the limit of the plotted sum is the
discounted-saving side of \eqref{eq:margin}. Because energy use rises with wealth in
levels, so does the saving. At mean wealth the saving never reaches $\psi_g$: the
typical household stays brown not because it discounts the future heavily but
because its energy bill is too small for the operating saving ever to repay the
cost. At about $25\,\bar a$ the breakeven arrives at ten quarters, half the durable's
expected life; at $40\,\bar a$ it arrives within seven quarters, well inside the durable's life,
and the switch clearly pays. The figure is an accounting cut: it holds the green
incumbent's demand fixed and abstracts from consumption smoothing and the taste
shocks that make the actual choice probabilistic. But it explains why the holding rate climbs from about $1.3\%$ to $84\%$ across
the wealth distribution (Figure~\ref{fig:ss_objects}), and it is the object the
crisis perturbs: a brown price shock scales up the saving flow and moves each
household's crossing left.

\begin{figure}[H]\centering
\figtitle{The private economics of one switch}
\IfFileExists{output/fig_adoption_payoff_import.pdf}{\includegraphics[width=.8\textwidth]{output/fig_adoption_payoff_import.pdf}}{\textit{[figure output/fig\_adoption\_payoff\_import.pdf: run run\_adoption\_payoff.py]}}
\label{fig:payoff}
\fignotes{Steady state, no shock. For the median discount-factor type at median
productivity and three wealth levels: cumulative energy saving from operating green,
discounted by interest and survival,
$S(T)=\sum_{t\le T}\bigl[(1-\delta_g)/(1+r)\bigr]^{t}
(\bar P^E_B-\bar P^E_G)\,\bar c^{\,G}_E(e,a)$, against the switching cost $\psi_g$
(dashed) and the expected durable life $1/\delta_g$ (dotted). Dots mark the
breakeven quarter.}
\end{figure}

Two further steady-state objects anticipate the policy results
(Figure~\ref{fig:ss_objects}). The switching \emph{flow} has a stock counterpart:
the green-holding rate rises from about $1.3\%$ at the constraint to $84\%$ at the top
of the wealth grid, so the durable is held disproportionately by the wealthy. The
marginal propensity to consume runs the other way, falling from about $0.5$ at the
constraint to near zero by $a\approx3.5\,\bar a$, with brown incumbents and green holders
close at each wealth level. These two gradients drive the incidence results of
Section~\ref{sec:flat}. A transfer to brown households reaches the high-MPC
constrained tail and stabilises consumption effectively; an energy-indexed transfer
does not, because the energy bill is flat as a share of spending
($3.8\%$ at the bottom of the wealth distribution against $3.5\%$ at the top) and so
rises with wealth in levels, sending the largest payments to the low-MPC wealthy.
Homothetic energy demand makes energy-indexed targeting regressive; the flat share
is the assumption behind that result, and the non-homothetic robustness check in
Section~\ref{sec:flat} relaxes it.

\begin{figure}[H]\centering
\figtitle{Green-durable holding and MPC across wealth}
\IfFileExists{output/fig_ss_objects_import.pdf}{\includegraphics[width=\textwidth]{output/fig_ss_objects_import.pdf}}{\textit{[figure output/fig\_ss\_objects\_import.pdf: run run\_ss\_objects.py]}}
\label{fig:ss_objects}
\fignotes{Aggregated over discount-factor types. Left: green-durable holding rate
against wealth (the stock behind the switching flow). Right: quarterly MPC against
wealth, by durable state. Wealth on a symmetric-log scale.}
\end{figure}

\subsection{The baseline wave under the two shocks}
\label{sec:baseline}

Figures~\ref{fig:baseline} and~\ref{fig:baseline_supply} report the baseline
responses to the two shocks, each with the adoption margin open (solid) against the
frozen counterfactual (dashed).

\begin{figure}[H]\centering
\figtitle{Price shock: baseline response}
\includegraphics[width=0.92\textwidth]{output/fig1_price_shock.png}
\label{fig:baseline}
\fignotes{Brown energy price shock, elastic supply, no policy. Adoption open (solid)
against the frozen-choice counterfactual at the common steady state (dashed).}
\end{figure}

\begin{figure}[H]\centering
\figtitle{Supply shock: baseline response}
\includegraphics[width=0.92\textwidth]{output/fig1_supply_shock.png}
\label{fig:baseline_supply}
\fignotes{Brown energy supply shock ($-10\%$ for 6 quarters, fixed quantity), no
policy. Adoption open (solid) against the frozen-choice counterfactual (dashed).}
\end{figure}

\subsection{The adoption channel}
\label{sec:signmap}

Under the price shock, the adoption channel's contribution
to cumulative output is $\Delta_{24}=-9.9$ (Table~\ref{tab:summary};
Figure~\ref{fig:decomposition} plots the margin's contribution across variables): the margin
is \emph{contractionary}. It follows a race between two resource flows in the
balance of payments (Section~\ref{sec:bop}). Adoption imports a switching cost $\psi_g D^{sw}_t$, a resource outflow. It is
paid once by each switcher, so it is front-loaded and scales with the number of
adopters. The energy saving $(\PEBpre-\PEGpre)\,C^{G}_{E,t}$ works the other way:
an inflow, booked as reduced imports, that accrues gradually over the life of the
green durable. Under the frozen counterfactual the front-loaded
cost dominates the windowed saving at every persistence. The force is the resource
cost of the crisis-time switching wave at the calibrated $\psi_g=0.77$
($\approx0.77\times C_{ss}$) rather than a general-equilibrium artefact: the
switching-cost outflow $\psi_g\times$ extra switchers cumulates to $+11.74$ on the
import side, and net of the energy-mix inflow it leaves a net import draw of
$+4.98$ (Figure~\ref{fig:adoption_flows}) that pulls resources out of domestic
absorption. Under the fixed-quantity (supply) closure a third, expansionary term
appears: adoption lowers fossil demand against a fixed supply, so the clearing
price falls and relaxes scarcity for \emph{all} households, which shifts
$\Delta_{24}$ up to $-0.82$. The marginal switcher's steady-state indifference,
$\psi_g\approx(\bar P^{E}_{B}-\bar P^{E}_{G})\,\bar C^{G}_{E}/(r+\delta_g)$, ties
the threshold to the cost ratio $\varrho$.

\paragraph{Booking caveat.} What the booking convention moves is the sign of the
adoption channels \emph{contribution} to output, $\Delta_{24}$, not the output
response itself: aggregate output falls under both shocks and every booking. At the
baseline $\alpha_F^{switch}=1$ the switching cost and green bill are pure imports, so
the contribution is negative and deepens the output fall; routing that spending onto
the home good raises $\Delta_{24}$, and under full domestic booking
($\alpha_F^{switch}=0$) the contribution turns positive, $+23.0$ under the price shock
and $+0.58$ under the supply shock. The policy ordering below is unaffected. We report $\alpha_F^{switch}=1$ throughout and
trace the full grid in Appendix~\ref{app:booking}. This dependence is a booking
convention rather than a result about the transition, and it is
why we phrase the paper's main result in terms of the \emph{margin's response} to
policy rather than its output sign.

\begin{figure}[H]\centering
\figtitle{Contribution of the adoption margin, by shock}
\includegraphics[width=\textwidth]{output/fig5_decomposition.png}
\label{fig:decomposition}
\fignotes{Full economy minus the frozen-choice counterfactual, by shock. Panels:
output, consumption, net exports. Price shock in black, supply shock in blue.}
\end{figure}

The race between the two resource flows is exact on the import side.
Figure~\ref{fig:adoption_flows} splits the adoption channel's contribution to
imports, full minus frozen, into the
front-loaded switching-cost outflow $\psi_g D^{sw}_t$ and everything else, the
latter dominated by the gradual energy saving from the brown-to-green mix shift.
The split is exact by construction. The switching cost dominates the saving at
every horizon under both shocks: under the price shock it cumulates to $+11.74$
against a $-6.76$ energy-mix inflow, a net import response of $+4.98$ that accounts
for the channel's contractionary output contribution ($\Delta_{24}=-9.9$). The
outflow peaks on impact, when household income is falling hardest, and this timing
is why booking it as an import makes the channel contractionary.

\begin{figure}[H]\centering
\figtitle{Adoption channel on the import side, by shock}
\includegraphics[width=\textwidth]{output/fig5b_adoption_flows.png}
\label{fig:adoption_flows}
\fignotes{Contribution of the adoption margin to imports, full economy minus the
frozen-choice counterfactual. The figure plots
the front-loaded switching-cost outflow $\psi_g D^{sw}$, the gradual energy-saving
inflow, and their sum. Level deviations $\times100$.}
\end{figure}

The adoption channel's output effect is a domestic-demand phenomenon. Decomposing
output one level below the domestic/export split (Figure~\ref{fig:output_channels}),
the frozen-choice counterfactual separates output into three exact pieces: exports,
domestic demand with the margin frozen (the AMRS real-income and substitution
channels), and the adoption channel's increment to domestic demand. Under the price
shock these cumulate to $y=-64.9 = (+7.4)_{\text{exports}} + (-62.9)_{\text{AMRS}} +
(-9.4)_{\text{adoption}}$. The adoption drag ($-9.4$) falls entirely on domestic
home-good demand: the switching-cost import outflow draws resources out of domestic
absorption. It is close in magnitude and opposite in sign to the export gain
($+7.4$), which is the AMRS competitiveness channel, so the two roughly cancel and
aggregate output tracks the AMRS domestic channel ($-62.9$).

\begin{figure}[H]\centering
\figtitle{Output response by channel}
\IfFileExists{output/fig_output_channels_import.pdf}{\includegraphics[width=0.8\textwidth]{output/fig_output_channels_import.pdf}}{\textit{[figure output/fig\_output\_channels\_import.pdf: run run\_output\_channels.py]}}
\label{fig:output_channels}
\fignotes{Price shock, no policy. Three exact channels: exports $c_H^{*}$, domestic
demand with the adoption margin frozen (AMRS), and the adoption channel's increment
to domestic demand. The three sum to $y$. Level deviations $\times100$.}
\end{figure}

The same response, read from the household budget rather than from output, is almost
entirely a general-equilibrium phenomenon. Decomposing consumption through the
household consumption Jacobians, $dC=\sum_x J^{C,x}\,dx$ over the household's inputs
and exact to first order, the direct energy-price effect is negligible ($\approx0$
cumulative, energy being a small budget share) and the collapse is the indirect
real-income channel ($-97.9$ of a $-109.3$ total, the real rate adding $-3.5$). The
energy shock reaches consumption by lowering factor income in general equilibrium,
not by raising the household's energy bill, the mechanism \citet{auclert2023energy} and the
broader HANK literature identify, here in an open economy with an adoption margin.

\subsection{The induced adopters}
\label{sec:induced}

The crisis draws in households who would not have
switched at the pre-crisis steady state. Figure~\ref{fig:induced_adopters}
characterises them with a comparative static: the switching probability of the median
type at a steady state with a permanently $10\%$ higher brown price, minus the
baseline field. The induced switching is positive throughout and rises with
productivity; at high productivity the induced margin reaches down to more moderate
wealth, pulling in households who sat just below the switching threshold. Weighted by
the steady-state distribution, the induced flow concentrates at moderate productivity
and low-to-moderate wealth, where the brown mass sits. The
construction is a comparative-static proxy: it uses a permanent price change rather
than the transitory crisis path, so it locates and signs the induced margin rather
than timing it.

\begin{figure}[H]\centering
\figtitle{Induced adopters}
\IfFileExists{output/fig_induced_adopters_import.pdf}{\includegraphics[width=0.8\textwidth]{output/fig_induced_adopters_import.pdf}}{\textit{[figure output/fig\_induced\_adopters\_import.pdf: run run\_induced\_adopters.py]}}
\label{fig:induced_adopters}
\fignotes{Comparative static: the extra quarterly switching probability induced by a
permanently $10\%$ higher brown price, median discount-factor type, over the
wealth--productivity grid. Contours mark the steady-state brown mass.}
\end{figure}
